\section{稳定性、正则与调度}

\subsection{Entropy 与多样性正则}

\begin{equation*}
\mathcal{L}_\pi = -\mathbb{E}_{s \sim \mathcal{D}, a \sim \pi} \left[\hat{A}(s,a)\log \pi(a|s)\right] + \beta \cdot \mathcal{H}(\pi(\cdot|s))
\end{equation*}

Entropy 项鼓励策略在更新阶段保留一定随机性，从而避免策略迅速收敛到次优动作。这在稀疏奖励或多模态动作空间尤为重要。

\begin{importantbox}{Entropy regularization 的实际量级}
在 PPO / SAC 中常见做法是把 entropy bonus $\beta$ 设为策略 loss 的 0.01~0.1，即保留 1%~10%  的 exploration energy。太小会陷入局部 optima，太大又会拖慢策略收敛。
\end{importantbox}

\subsection{Trust region 与 clipped objective}

PPO 的 clipped objective：
$${L}^{CLIP} = \mathbb{E} \left[ \min\left( r_t(\theta)\hat{A}_t, \; \text{clip}(r_t(\theta), 1-\epsilon, 1+\epsilon)\hat{A}_t \right) \right]$$
其中 $r_t(\theta) = \frac{\pi_\theta(a_t|s_t)}{\pi_{\theta_{\text{old}}}(a_t|s_t)}$。

把 KL 约束或者 clip 机制当作 trust region，可以防止策略步幅过大。讲者提示：在 PPO 里监控 $KL(\pi_{\theta_{\text{old}}}||\pi_\theta)$，一旦跳过阈值就退回到上一次 checkpoint。

\begin{knowledgebox}{Trust region vs clipping}
TRPO 通过精确算梯度约束，计算量大；PPO 则把约束用 clip 近似，计算简单且易于并行，是当前主流首选算法。
\end{knowledgebox}

\subsection{调度与 target network}

Critic 的 Target Network（如 SAC 的慢速权重）通过 Polyak averaging 让目标值平滑。例如：
$$\theta_{\text{target}} \leftarrow \tau \theta + (1 - \tau) \theta_{\text{target}}$$

较小的 $\tau$ 值让 target network 更新缓慢，提高稳定性，但也可能在非平稳环境里反应迟钝。

\begin{warningbox}{Target network 更新太慢会跟不上策略}
如果 $\tau$ 设得过小，critic 会一直追随已经“过时”的 target，导致引导 Actor 估计过于乐观，对新样本适应慢。通常把 $\tau$ 设在 0.005~0.05 之间。
\end{warningbox}

\subsection{本章小结}

稳定性技巧（entropy、clip objective、target network）是 Actor-Critic 能否快速收敛的关键。监控 KL 与 advantage 的分布可以及时发现出界更新。

